\documentclass[preview]{standalone}
\usepackage{amsmath}
\usepackage{amssymb}
\usepackage{stellar}
\usepackage{definitions}
\usepackage{bettelini}
\begin{document}
\id{diffeq-separable}
\genpage
\section{Separable equations}

\begin{snippetdefinition}{separale-equation-definition}{Separable equation}
    A first-order equation is \emph{separable} if it has the form
    \(y'=h(x)u(y)\), where \(h\colon I\fromto\realnumbers\) and
    \(u\colon J\fromto\realnumbers\) are continuous on open intervals.
    A related separated form is \(a(y)y'=h(x)\); wherever \(a(y)\neq0\),
    this is the normal form \(y'=h(x)/a(y)\).
\end{snippetdefinition}

\begin{snippetproposition}{diffeq-separable-equilibria}{Constant solutions}
    Let \(h\colon I\fromto\realnumbers\) be continuous and
    \(u\in\continuityclass^1(J)\), with \(I,J\) open intervals.
    For every \(y_0\in J\) with \(u(y_0)=0\), the constant \function \(y=y_0\)
    solves \(y'=h(x)u(y)\) on \(I\), and is the unique solution with
    that initial value. If \(h\not\equiv0\), these are all constant solutions on \(I\).
\end{snippetproposition}

\begin{snippetproof}{diffeq-separable-equilibria-proof}{diffeq-separable-equilibria}{Constant solutions}
    Substitute \(y=y_0\) into the equation. Local uniqueness applies because
    \(f(x,y)=h(x)u(y)\) and \(\partial_y f=h(x)u'(y)\) are continuous.
    If \(h(x_1)\neq0\), the equation for a constant solution implies \(u(y_0)=0\).
\end{snippetproof}

\begin{snippettheorem}{separale-equation-sol}{Separation of variables}
    Let \(I,J\) be open intervals, \(h\in C(I)\), \(u\in C^1(J)\),
    \(x_0\in I\), and \(y_0\in J\) with \(u(y_0)\neq0\).
    On the component \(J_0\) of \(\{s\in J:u(s)\neq0\}\) containing \(y_0\), set
    \[
        F(v)=\integral[y_0][v][\frac{1}{u(s)}][s],
        \qquad G(x)=\integral[x_0][x][h(t)][t].
    \]
    The initial value problem \(y'=h(x)u(y)\), \(y(x_0)=y_0\),
    has the local solution \(y=F^{-1}\circ G\).
    Equivalently, its solution in \(J_0\) is characterized by
    \[
        \integral[y_0][y(x)][\frac{1}{u(s)}][s]
        =\integral[x_0][x][h(t)][t].
    \]
    More generally, if \(H'=a\) and \(G'=h\), a differentiable \function
    \(y\) solves \(a(y)y'=h(x)\) \ifandonlyif \(H(y(x))=G(x)+C\).
\end{snippettheorem}

\begin{snippetproof}{separale-equation-sol-proof}{separale-equation-sol}{Separation of variables}
    Near \(x_0\), continuity gives \(u(y(x))\neq0\). Divide by this quantity
    and integrate:
    \[
        \integral[x_0][x][\frac{y'(t)}{u(y(t))}][t]
        =\integral[x_0][x][h(t)][t].
    \]
    The chain rule, or the substitution \(s=y(t)\), identifies the left side
    with \(F(y(x))\). Since \(F'=1/u\) is continuous and never zero on \(J_0\),
    \(F\) is strictly monotone and has a differentiable inverse.
    As \(G(x_0)=F(y_0)=0\), the inverse expression is defined near \(x_0\).
    Differentiating it proves the converse. For the separated form,
    \((H\circ y)'=a(y)y'\), so integration and differentiation prove both implications.
\end{snippetproof}

\begin{snippetnote}{diffeq-separation-domain-warning}{Domains and lost solutions}
    Dividing \(y'=h(x)u(y)\) by \(u(y)\) discards the constant solutions
    at zeros of \(u\); list them first. With primitives \(F'=1/u\), \(G'=h\),
    the relation is \(F(y)=G(x)+C\). Inverting requires choosing a branch
    on which \(u\neq0\) and restricting \(G(x)+C\) to the range of \(F\).
    The component containing the initial point determines the relevant interval.
\end{snippetnote}

\section{Examples}

\begin{snippetexample}{diffeq-separable-quartic-example}{A global positive solution}
    Solve \(y'=4x^3/y\), \(y(0)=2\), in
    \(\Omega=\realnumbers\times(0,+\infty)\).
    Here \(h(x)=4x^3\), \(u(y)=1/y\), so there are no constant solutions.
    Separation gives
    \[
        yy'=4x^3,\qquad (y^2/2)'=(x^4)',\qquad y^2/2=x^4+C.
    \]
    The initial value gives \(C=2\). The positive branch is therefore
    \[
        y(x)=\sqrt{2x^4+4},\qquad x\in\realnumbers.
    \]
    The radicand is positive everywhere, so this solution is global.
\end{snippetexample}

\begin{snippetexample}{diffeq-separable-exponential-example}{Dependence of the domain on the initial value}
    Consider \(y'=e^{x-y}\), \(y(0)=a\), with \(a\in\realnumbers\).
    Here \(h(x)=e^x\), \(u(y)=e^{-y}\), and the equation is defined on
    \(\realnumbers^2\), with no constant solutions. Separation yields
    \[
        e^y y'=e^x,\qquad (e^y)'=(e^x)',\qquad e^y=e^x+C.
    \]
    At zero, \(e^a=1+C\), giving
    \[
        y(x)=\ln(e^x+e^a-1).
    \]
    Its domain requires \(e^x>1-e^a\).
    If \(a\geq0\), the right-hand side is nonpositive and the maximal interval
    is \(\realnumbers\). If \(a<0\), it is
    \[
        I=(\ln(1-e^a),+\infty).
    \]
    In the latter case the finite left endpoint is a vertical asymptote.
\end{snippetexample}

\begin{snippetexample}{diffeq-separable-implicit-example}{An implicit solution}
    Consider \(y'=2x/(ye^y)\), \(y(1)=3\), on the positive \(y\)-half-plane.
    Here \(h(x)=2x\) and \(u(y)=1/(ye^y)\). Multiplication and integration give
    \begin{align*}
        ye^y y'&=2x,\\
        \integral[1][x][y(t)e^{y(t)}y'(t)][t]
        &=\integral[1][x][2t][t],\\
        \integral[3][y(x)][se^s][s]&=x^2-1.
    \end{align*}
    Integration by parts gives \(\int se^s\,ds=(s-1)e^s\), hence
    \[
        (y(x)-1)e^{y(x)}-2e^3=x^2-1.
    \]
    This cannot be solved for \(y\) using elementary \function[functions].
    Nevertheless \(v\mapsto(v-1)e^v\) has derivative \(ve^v>0\) for \(v>0\);
    its range there is \((-1,+\infty)\). Since \(x^2-1+2e^3>-1\),
    the implicit formula determines a positive solution for every real \(x\).
\end{snippetexample}
\end{document}
